\documentclass[11pt,a4paper]{article}
\usepackage[brazil]{babel}
\usepackage[utf8]{inputenc}
\usepackage[T1]{fontenc}
\usepackage{amsmath,amssymb}
\usepackage[margin=2cm]{geometry}
\usepackage[hidelinks]{hyperref}

\title{Fichamento: \cite{Finazzi2009}}
\author{Tiago Rodrigues}
\date{\today}

\begin{document}
\maketitle

\section*{Pergunta central}
% Que problema o artigo tenta resolver?

\section*{Ideia principal}
% Em 3 a 5 frases, com suas palavras.

\section*{Equações-chave}
% A métrica, o tensor energia-momento, a desigualdade... e o significado de cada símbolo.
\begin{equation*}
\end{equation*}

\section*{Resultados}

\section*{Limitações e críticas}
% O que os artigos seguintes apontaram sobre este?

\section*{Conexões com outros artigos}

\section*{Dúvidas}

\bibliographystyle{unsrt}
\bibliography{../../referencias}
\end{document}
